\newpage
\section{Ejemplo: ascenso con inclinación lateral}
Consideremos dos instantes separados por $\Delta t=0.10\,\mathrm{s}$ durante el ascenso. El cohete realiza una rotación suave y controlada hacia el lateral $+x$ del marco local. La inclinación de su eje longitudinal pasa de $10^\circ$ a $20^\circ$ respecto a la vertical, mediante una rotación alrededor del eje $y$. El carácter controlado de esta maniobra es una hipótesis del ejemplo, no una función de control definida aquí.

Los valores son ilustrativos. Suponemos lecturas sincronizadas y calibradas, orientación conocida por una estimación independiente apoyada en el giróscopo, y un acelerómetro próximo al centro de masa, de modo que se desprecia el efecto de su desplazamiento respecto a ese punto. Las componentes indicadas son fuerza específica en los ejes del sensor, expresada en $\mathrm{m/s^2}$.
\begin{center}
\begin{tabular}{c|c|c|c}
Instante & Presión & Inclinación & Lectura $(f_x,f_y,f_z)$ \\
\hline
$t_1$ & $100000\,\mathrm{Pa}$ & $10^\circ$ & $(1.00,\ 0.00,\ 15.00)$ \\
$t_2=t_1+0.10\,\mathrm{s}$ & $99990.5\,\mathrm{Pa}$ & $20^\circ$ & $(2.00,\ 0.00,\ 14.00)$
\end{tabular}
\end{center}
El movimiento se limita al plano $xz$ para facilitar el cálculo; por eso $f_y=0$. Tomamos $\bar T=293.15\,\mathrm{K}$ y una velocidad previamente estimada $V_r(t_1)=8.00\,\mathrm{m/s}$. Esta condición inicial es necesaria: dos lecturas de aceleración solo proporcionan el cambio de velocidad.

\subsection{Aporte del barómetro}
Con las dos presiones, la velocidad barométrica resulta
\begin{equation}
 V_b=\frac{287(293.15)}{9.81(0.10)}
 \ln\left(\frac{100000}{99990.5}\right)
 \approx8.148\,\mathrm{m/s}.
\end{equation}
La caída de presión equivale a un incremento de altura aproximado de $0.815\,\mathrm{m}$. Esta velocidad corresponde al promedio del intervalo y es positiva porque el cohete asciende.

\subsection{Proyección de las componentes del acelerómetro}
Para la rotación lateral definida, la transformación del marco del sensor al marco local es
\begin{equation}
 \mathbf R_{NB}(\theta)=
 \begin{bmatrix}
 \cos\theta&0&\sin\theta\\
 0&1&0\\
 -\sin\theta&0&\cos\theta
 \end{bmatrix}.
\end{equation}
Por tanto, las aceleraciones de trayectoria lateral y vertical son
\begin{equation}
 a_x=f_x\cos\theta+f_z\sin\theta,\qquad
 a_z=-f_x\sin\theta+f_z\cos\theta-g.
\end{equation}
La componente $f_z$ del sensor no es la componente vertical local cuando el cuerpo está inclinado. Además, la inclinación del cuerpo y la dirección de la aceleración no tienen por qué coincidir.

\newpage
En el primer instante se obtiene
\begin{align}
 a_{z,1}&=-1.00\sin10^\circ+15.00\cos10^\circ-9.81
 \approx4.788\,\mathrm{m/s^2},\\
 a_{x,1}&=1.00\cos10^\circ+15.00\sin10^\circ
 \approx3.590\,\mathrm{m/s^2}.
\end{align}
En el segundo instante,
\begin{align}
 a_{z,2}&=-2.00\sin20^\circ+14.00\cos20^\circ-9.81
 \approx2.662\,\mathrm{m/s^2},\\
 a_{x,2}&=2.00\cos20^\circ+14.00\sin20^\circ
 \approx6.668\,\mathrm{m/s^2}.
\end{align}
Los vectores de aceleración en el marco local son, aproximadamente,
\begin{equation}
 \mathbf a_{N,1}=\begin{bmatrix}3.590\\0\\4.788\end{bmatrix}\,\mathrm{m/s^2},
 \qquad
 \mathbf a_{N,2}=\begin{bmatrix}6.668\\0\\2.662\end{bmatrix}\,\mathrm{m/s^2}.
\end{equation}
Existe aceleración lateral, pero solo $a_z$ interviene en la actualización de velocidad vertical. Si se restara directamente $g$ a $f_{z,2}$, se obtendría $4.19\,\mathrm{m/s^2}$ en lugar de $2.662\,\mathrm{m/s^2}$: se estaría ignorando la orientación.

\subsection{Predicción inercial y velocidad combinada}
Integrando las componentes verticales con la regla del trapecio,
\begin{equation}
 V_a(t_2)=8.00+\frac{4.788+2.662}{2}(0.10)
 \approx8.373\,\mathrm{m/s}.
\end{equation}
Para ilustrar la combinación elegimos $\tau=0.40\,\mathrm{s}$, de modo que
\begin{equation}
 \alpha=\frac{0.40}{0.40+0.10}=0.80.
\end{equation}
Esta elección es didáctica; no representa un ajuste validado. Con los valores calculados antes de redondear,
\begin{equation}
 V_r(t_2)=0.80\,V_a(t_2)+0.20\,V_b
 \approx\boxed{8.328\,\mathrm{m/s}}.
\end{equation}
El resultado indica que el cohete sigue ascendiendo a una velocidad vertical estimada de aproximadamente $8.33\,\mathrm{m/s}$, aunque su aceleración también tiene una componente lateral. No es la magnitud de la velocidad total del cohete.

En este cálculo, la predicción inercial corresponde a $t_2$, mientras que el barómetro aporta un promedio entre $t_1$ y $t_2$. La combinación aproxima la velocidad actual bajo la hipótesis de un intervalo corto; esa diferencia temporal introduce error durante las aceleraciones. Los decimales mostrados permiten seguir la aritmética, pero no implican esa precisión física en los sensores.
\end{document}
